\section{Diagnostic Methodology}\label{sec:method}
Our method is a diagnosis: research questions, theory that fixes what sharing can and cannot represent, one main hypothesis with testable predictions, and tests for the alternatives. The \crs{} is one of these tests.

\subsection{Research questions}
\begin{description}
\item[RQ1] Is the Full$\to$shared accuracy drop of DirectAU a loss of what magnitude contributes, of what direction contributes, or both?
\item[RQ2] Under sharing, does magnitude follow buckets instead of items?
\item[RQ3] Do alternative explanations account for the magnitude part of the drop?
\item[RQ4] Does restoring item counts let DirectAU outperform BPR again under the same table, and is the gain item-specific?
\item[RQ5] Does the pattern hold for other sharing schemes, bucket resolutions and LightGCN?
\end{description}
Costs and robustness are treated in \S\ref{sec:tradeoffs}.

\subsection{What a shared table can represent}\label{sec:theory}
Each result states its assumptions in its first clause; proofs are in Appendix~\ref{app:proofs}.

\begin{lemma}[User scale]\label{lem:user}
Under raw, cosine and count-restored scoring, rescaling $\mathbf{e}_u$ by any $a>0$ leaves the ranking of user $u$ unchanged. Under raw${}+\beta p$ it does not.
\end{lemma}
Under the first three scorings, popularity can therefore reach a ranking only through item-side quantities: magnitude, direction, or an explicit item term.

\begin{lemma}[Intra-bucket score ties]\label{lem:ties}
Under single-row sharing in MF, items with $b(i)=b(j)$ receive identical scores for every user under any scoring that depends on the item only through $\mathbf{q}_i$.
\end{lemma}
The lemma does not cover two-row schemes unless both rows coincide, nor LightGCN, whose propagated item embeddings differ within a bucket. Any item term, such as $\beta p_i$, breaks these intra-bucket score ties. They are common: on Gowalla, 87\% of evaluated users have a tie that straddles rank 20, and the mean tied group at that boundary has 14.7 items.

\begin{proposition}[Scale invariance]\label{prop:scale}
$\mathcal{L}_{\mathrm{DAU}}$ depends on each item embedding only through $\mathbf{q}_i/\|\mathbf{q}_i\|$.
\end{proposition}
\begin{corollary}[A learned per-item scale receives no gradient]\label{cor:scale}
Let $\mathbf{q}_i=w_i\,\mathbf{E}_{b(i)}$ with a learned scalar $w_i>0$, the form used by MEmCom~\cite{pansare2022memcom}. Then $\partial\mathcal{L}_{\mathrm{DAU}}/\partial w_i=0$, and $w_i$ does not affect cosine or count-restored scores.
\end{corollary}
The corollary is limited to the DirectAU loss: a weight-decay term on $w_i$ still shrinks it, and BPR is not scale-invariant.

\begin{remark}[Norm dynamics]
Under plain SGD without weight decay, Proposition~\ref{prop:scale} and Euler's identity give $\langle\nabla_{\mathbf{q}}\mathcal{L}_{\mathrm{DAU}},\mathbf{q}\rangle=0$, so $\|\mathbf{q}_{k+1}\|^2=\|\mathbf{q}_k\|^2+\eta^2\|\mathbf{g}_k\|^2$ and magnitudes cannot shrink. The same holds for a shared row, whose gradient is a sum of terms orthogonal to it. Adam's per-coordinate scaling and weight decay break this argument, so we report magnitude growth under our optimiser as an observation (\S\ref{sec:rq2}), not as a theorem.
\end{remark}

\begin{proposition}[Correlation-ratio bound]\label{prop:bound}
For any item-level target $t$ with $\Var(t)>0$ and any bucket-level quantity $f_i=g(b(i))$ with $\Var(f)>0$, $\corr(f,t)^2\le\eta^2(t)$, with correlations over items with uniform weight. Equality holds for $g=\mu_b(t)$.
\end{proposition}
The proof is one line: write $t=\mu_{b}(t)+\varepsilon$; then $\sum_{j:b(j)=b}\varepsilon_j=0$ for every bucket, so $\varepsilon$ is orthogonal to every bucket-level quantity and $\mathrm{Cov}(f,t)=\mathrm{Cov}(f,\mu_b(t))$; Cauchy--Schwarz finishes it. Two readings follow.
With $t=p$, no bucket-level quantity, including a shared row's magnitude, a learned bucket bias or a bucket-level prior, can explain more than $\eta^2_p=0.36$, $0.48$ and $0.36$ of log-count variance under GraphHash on Gowalla, Yelp2018 and AmazonBook.
With $t=\mathrm{rank}(c)$, a bucket-level magnitude has bucket-level average ranks, so $\rho_{\mathrm{item}}^2\le\eta^2(\mathrm{rank}\,c)=0.271$, $0.437$ and $0.303$. The bound is sufficient, not necessary: under double-frequency hashing on Gowalla $\eta^2_p=0.999$, yet $\rho_{\mathrm{item}}$ is only $0.03$.
We call the item-specific part $\varepsilon_i=p_i-\mu_{b(i)}(p)$ the \emph{within-bucket residual}. By Lemma~\ref{lem:ties} and Proposition~\ref{prop:bound}, a single-row shared table cannot store it.

\subsection{Main hypothesis and its predictions}\label{sec:hypothesis}
\textbf{H\textsubscript{pool} (\pooling).} Under sharing, a row's magnitude is set by its bucket's traffic, so magnitude stops carrying item popularity, and DirectAU, which on full tables relies on magnitude at scoring time, loses that contribution.
H\textsubscript{pool} predicts:
(P1)~$\rho_{\mathrm{bucket}}\gg\rho_{\mathrm{item}}$ under sharing;
(P2)~for DirectAU, cosine $\approx$ raw under sharing but cosine $\ll$ raw on full tables;
(P3)~restoring $p_i$ at scoring time recovers a gain carried mostly by $\varepsilon$;
(P4)~datasets where magnitude contributes more on the full table lose more under sharing.

\subsection{Separating magnitude from direction}\label{sec:decomposition}
Let $R(s,m)$ be Recall@20 of a table $s\in\{\mathrm{full},\mathrm{shared}\}$ under scoring $m\in\{\mathrm{raw},\cos\}$, all measured on the same users. Define the value of magnitude on each table, $M_s=R(s,\mathrm{raw})-R(s,\cos)$, and the drop under direction-only scoring, $D=R(\mathrm{full},\cos)-R(\mathrm{shared},\cos)$. Then, exactly,
\begin{equation}\label{eq:decomp}
\underbrace{R(\mathrm{full},\mathrm{raw})-R(\mathrm{shared},\mathrm{raw})}_{\text{raw drop}}
= D + \underbrace{(M_{\mathrm{full}}-M_{\mathrm{shared}})}_{I}.
\end{equation}
$I$ is the interaction between sharing and magnitude, a difference in differences: the part of the drop that exists only because sharing removes what magnitude contributed. We call $M_{\mathrm{full}}$ the magnitude channel and $D$ the direction channel. Equation~\eqref{eq:decomp} is an exact identity on measured recalls, and because all four cells are evaluated on the same users, every term carries a paired user-bootstrap interval.

\subsection{Alternative explanations}\label{sec:alternatives}
Table~\ref{tab:alternatives} lists nine alternatives to H\textsubscript{pool}, each with a targeted test: controlled scoring or training contrasts for A1--A4 and A6--A8, the stopping protocol for A5, and a proof for A9. Results are in \S\ref{sec:rq3}.

\begin{table}[t]
\caption{Alternatives to H\textsubscript{pool} and their tests.}\label{tab:alternatives}
\centering\small
\begin{tabular}{@{}p{0.40\columnwidth}p{0.54\columnwidth}@{}}
\toprule
Alternative & Test\\
\midrule
A1 Lost capacity explains all of it & Full table trained with equal item magnitudes (all parameters kept)\\
A2 Ties alone & Break ties by popularity, without the prior\\
A3 Bucket-level popularity suffices & Fixed-$\beta$ $\mu/\varepsilon$ decomposition; shuffle counts within buckets\\
A4 Magnitude carries other information & Full table: replace magnitude by a fitted function of $p$; global and within-decile shuffles\\
A5 Undertraining & Protocol: only runs stopped by patience\\
A6 Initialisation scale & Init-matched reruns\\
A7 Weight decay & Training with weight decay 0\\
A8 Generic popularity & Popularity-only ranking\\
A9 A learned scale would suffice & Corollary~\ref{cor:scale}\\
\bottomrule
\end{tabular}
\end{table}

\subsection{The intervention as a test}\label{sec:intervention}
If H\textsubscript{pool} holds, replacing pooled magnitude with item counts should let DirectAU outperform BPR again \emph{under the same shared table}. We test this with the \crs{} for DirectAU, which matches its normalised training geometry, and with raw${}+\beta p$ for BPR, which matches its dot-product training. Both pairings were registered before the final runs. Only $\beta$ is selected, on validation Recall@20, from $\{0,0.01,0.02,0.05,0.1,0.2,0.5,1,2,5\}$, and the test set is scored once. The conclusion does not depend on the pairing: choosing each loss's correction by the same validation rule gives DirectAU the win on all nine dataset--seed pairs, and on Yelp2018 and AmazonBook DirectAU's weaker correction still beats BPR's stronger one (\S\ref{sec:robustness}). All modes are reported in Appendix~C.

To separate bucket-level from item-specific popularity (A3), we split the prior at the selected $\beta^\ast$ as $\beta^\ast p_i=\beta^\ast\mu_{b(i)}+\beta^\ast\varepsilon_i$ and score the four cells base, $+\beta^\ast\mu$, $+\beta^\ast\varepsilon$ and $+\beta^\ast p$. Scores decompose exactly but Recall does not, so we report main effects, their interaction, and two-factor Shapley shares of the \emph{correction gain}. All cells rank with exact intra-bucket ties broken by a random key, never by adding a tiny $\beta p$.
